%% This file is an example of using beamerthemeedmaths.sty to style a Beamer
%% presentation. While the style is LPPL 1.3 licensed, this example is given
%% under a 0BSD license (https://choosealicense.com/licenses/0bsd), meaning
%% it can be freely used, copied, modified, and/or redistributed with or
%% without attribution. For instructions on using the style, see the README.md
%% file or https://github.com/UoE-School-of-Mathematics/LaTeX-Templates.

\documentclass[notes]{beamer}
\usetheme{edmaths}
\usepackage{lipsum}
\usepackage[overridenote=true]{pdfpc}

\pdfcompresslevel=9
\pdfobjcompresslevel=3

\title[Insert Title (short version of title)]{Title of Presentation}
\subtitle{Insert Subtitle (Optional)}  % Optional
\author{Author Name}
\date{n\textsuperscript{th} Month YYYY}

\begin{document}


\begin{frame}
    \titlepage{}
\end{frame}
\note{
    This generates notes for pdfpc. \\
    \\
    These notes also appear
    on the handout/article versions.
}


\begin{frame}[t]{Header title bold}\label{sec:Slide1}
    Main body text Sans. Here is a way to \alert{highlight text}.
    \begin{itemize}
    \item Bullet points
    \item Sub-lists\ldots
        \begin{itemize}
            \item subitem
        \end{itemize}
    \item \lipsum[2]
    \end{itemize}
\end{frame}
\note{
    This is a note for the first slide. Mention how something happens.
}


\begin{frame}{Inserting a figure, some maths\ldots}
    Example of a reference to another part of the presentation (slide~\ref{sec:Slide1}). Also an example of a slide not vertically aligned to top.\newline

    \begin{itemize}
        \item Here's a figure:
            \begin{figure}[h]
                \centering
                \colorbox{blue}{\includegraphics[width=.13\linewidth]{images/edinburgh-logo}}
                \caption{The University of Edinburgh's logo}\label{fig:edinburgh-logo}
            \end{figure}
        \item And here's a famous equation:
            \begin{eqnarray}
                \label{eq:schroedinger}
                i\hbar {\frac{\partial}{\partial t}}|\psi (t)\rangle ={\hat{H}}|\psi (t)\rangle
            \end{eqnarray}
    \end{itemize}
\end{frame}


\begin{frame}[t]{References}
    The stylesheet {\bf\texttt{edmaths.sty}} provides a styling for LaTeX{} documents which satisfies University of Edinburgh requirements~\cite{koeppe2007}. \newline        

    \lipsum[4]\newline

    References can be added to the bottom of slides, or on a new slide.
    \vfill 
    \bibliography{example-references}
    \bibliographystyle{ieeetr}
\end{frame}


\end{document}
